\section{Giới thiệu}
\subsection{Đặt vấn đề}

Thành phố Hồ Chí Minh hiện đang đối mặt với thách thức nghiêm trọng về ùn tắc giao thông, đặc biệt trong khung giờ cao điểm buổi sáng (6–9 giờ) và buổi chiều (16–19 giờ). Sự gia tăng nhanh chóng về số lượng xe cá nhân – trong đó xe máy chiếm phần lớn tổng phương tiện – đã vượt xa năng lực tiếp nhận của mạng lưới hạ tầng hiện hữu. Kết quả là thời gian di chuyển trung bình cho quãng đường chỉ 10 km có thể kéo dài lên đến 60–90 phút, gây lãng phí nhiên liệu, tăng chi phí vận hành và làm giảm năng suất lao động toàn xã hội. Bên cạnh đó, lượng khí thải thoát ra trong điều kiện lưu thông ì ạch làm suy giảm chất lượng không khí, đe dọa sức khỏe cộng đồng và gia tăng nguy cơ mắc các bệnh đường hô hấp. Trong bối cảnh này, việc xây dựng một hệ thống phân tích dữ liệu lớn có khả năng thu thập và xử lý hình ảnh từ hàng ngàn camera giao thông công cộng, cung cấp bản đồ nhiệt mật độ phương tiện theo thời gian thực và gợi ý lộ trình tối ưu cho người tham gia giao thông không chỉ là giải pháp kỹ thuật mà còn là yêu cầu cấp thiết nhằm giảm thiểu ùn tắc, tiết kiệm nguồn lực và cải thiện chất lượng cuộc sống đô thị.
\subsection{Các hướng giải quyết liên quan}

Các giải pháp truyền thống như đếm xe bằng cảm biến cố định (vòng từ, hồng ngoại) hay phân tích tín hiệu GPS trên phương tiện công cộng đã chứng minh hiệu quả nhất định, nhưng vẫn vướng hạn chế về chi phí lắp đặt, khả năng mở rộng và độ bao phủ dữ liệu. Trong khi đó, các nền tảng như Google Maps hoặc Here Maps chủ yếu khai thác dữ liệu crowdsourcing và GPS từ người dùng, chưa tận dụng triệt để mạng lưới camera giao thông sẵn có. Xu hướng hiện nay chuyển sang kết hợp Computer Vision và AI trên dòng dữ liệu hình ảnh: sử dụng các mô hình học sâu (YOLO, Faster R‑CNN…) để phát hiện/đếm phương tiện trong video, sau đó áp dụng các mô hình chuỗi thời gian (LSTM, Prophet) hoặc mô hình đồ thị (Graph Neural Networks) để dự báo mật độ và nguy cơ ùn tắc trong ngắn hạn. Kết quả được quy đổi theo ô lưới (grid) trên bản đồ và hiển thị dưới dạng heatmap trực quan; đồng thời tích hợp vào ứng dụng di động để gợi ý lộ trình tránh “điểm nóng”

HAHHAHAHAHAHAH